Equilibrium statistical mechanics is by now a mature field. One of its most remarkable results is the universal structure of the equilibrium distribution: any system that can exchange heat with a (single) large enough thermal bath eventually relaxes into an equilibrium distribution, the Boltzmann distribution, that is completely determined by the system's Hamiltonian and the bath's temperature. At the microscopic level, the stability of this distribution is guaranteed by the principle of detailed balance, which ensures the absence of net probability currents \cite{seifert2012stochastic}.
% where the probability to find the system in a microstate $\mu_i$ is given by the Boltzmann distribution, $P(\mu_i) = e^{-H(\mu_i)/k_BT}/Z$, where $H(\mu_i)$ is the value of the Hamiltonian---the energy functional of the system---at this microstate, $T$ is the bath temperature, $k_B$ is the Boltzmann constant and $Z$ is the partition function, which is the corresponding normalization.
% Unfortunately, we do not know of an analog result for systems that are not in thermal equilibrium. In some nonequilibrium scenarios, e.g. in periodically driven system or glassy dynamics, the system does not relax into a steady state, and the probability to find the system in a given microstate changes over time or strongly depends on the exact initial condition. But even in the simplest type of nonequilibrium, where detailed balance is violated but yet the system relaxes into a NESS (nonequilibrium steady state), we have no simple formula for this steady state, and finding it is a difficult problem that to the best of our knowledge has to be solved case by case.

Unfortunately, no such universal principle exists for systems that are not in thermal equilibrium. Some nonequilibrium scenarios, such as glassy or periodically driven systems, may never reach a stationary state. But even the seemingly simpler case of a Non-Equilibrium Steady State (NESS) poses a significant theoretical challenge \cite{sasa2006steady}. In a NESS, the violation of detailed balance is manifested in persistent probability fluxes \cite{seifert2012stochastic, sasa2006steady, zia2007probability}, such that the stationary distribution alone does not uniquely describe the system \cite{zia2007probability}. As a consequence, there is no direct analog of the Boltzmann distribution for a NESS, and exact solutions typically require a case-by-case analytical approach.

There are many physical scenarios where the detailed balance condition does not hold and the system is driven out of equilibrium: some degrees of freedom (DoF) might be driven microscopically as in active matter \cite{marchetti2013hydrodynamics, bechinger2016active}; feedback control protocols \cite{seifert2012stochastic, sagawa2012nonequilibrium, parrondo2015thermodynamics, goerlich2025experimental} and time dependent driving \cite{jung1993periodically, schmiedl2008efficiency, raz2016mimicking,busiello2018similarities, koyuk2020thermodynamic}; a system connected to two or more baths with different temperatures, chemical potentials or voltages \cite{zia2007probability, derrida2007non}; or the system can be influenced by some external non-conservative force \cite{seifert2012stochastic, sasa2006steady, chernyak2006path}. The two latter cases are the main focus of this manuscript. 

The study of systems coupled to multiple reservoirs originates in the very foundation of classical thermodynamics, where the heat engine operating between a hot and a cold reservoir serves as the paradigmatic setup studied by Carnot \cite{carnot1872reflexions}. Clausius's formulation of the second law of thermodynamics also concerns heat flows between reservoirs \cite{clausius1879mechanical}. Later on, the transport associated with gradients in the value of thermodynamic quantities was considered by Onsager in his seminal works on reciprocal relations \cite{onsager1931reciprocal}. Systems coupled to multiple baths remain of great interest to date, where typically a microscopic model is considered instead of the earlier macroscopic descriptions. These include boundary-driven models such as asymmetric simple exclusion processes (ASEP) \cite{krug1991boundary, derrida1998exactly, derrida2007non} or oscillator chains \cite{kipnis1982heat, eckmann1999non}. More recently, there has also been great interest systems where different DoF are coupled to a different heat bath within the framework of stochastic thermodynamics \cite{seifert2012stochastic}. Notable examples are Brownian motors and refrigerators \cite{van2004microscopic, van2006brownian, filliger2007brownian}, which are a modern take on the classical macroscopic heat engine. These ideas have been generalized and studied extensively both theoretically and experimentally \cite{ dotsenko2013two, mancois2018two, bae2021inertial, baldassarri2020engineered,cerasoli2018asymmetry, cerasoli2022spectral, argun2017experimental, martins2026role, chang2021autonomous, dos2021stationary, pietzonka2018universal, lin2022stochastic, movilla2021energy, abdoli2022tunable, chiang2017electrical, ciliberto2013heat, PhysRevE.92.032118, abdoli2025quadrupolar, visco2006work, berut2014energy, berut2016stationary, fogedby2011bound, fogedby2017minimal, leighton2024information, netz2020approach} in the past two decades. In these systems, the simultaneous coupling to several reservoirs imposes incompatible equilibrium conditions; consequently, the system cannot relax to a standard Boltzmann distribution, but settles instead into a NESS characterized by persistent currents and non-vanishing entropy production. Given the well-established nature of these systems, it is particularly compelling to determine when seemingly distinct driving mechanisms, such as the non-conservative forces discussed below, can be formally mapped onto this reservoir-driven paradigm for nonequilibrium thermodynamics.

In principle, all the fundamental forces in nature are conservative, and can be derived from a potential energy function. However, in cases where a complete microscopic description is unavailable for some parts of the total system, it  is often necessary to approximate the dynamic through non-conservative forces. For example, drag forces are often introduced to account for the influence of a medium on the motion of a particle. Generally, non-conservative forces can be \emph{dissipative}, which means that the dynamics of the system are not reversible and do not conserve phase space volume. Such non-conservative forces do not necessarily violate the detailed balance condition (and evidently, they are necessary for the equilibration of an excited system). However, there also exist reversible non-conservative forces, which conserve phase space volume but are not derived from a potential. Specifically, we are interested in the case of force fields $\mathbf{F}(\mathbf{x})$ that depend only on the position $\mathbf{x}$ of the particles but not on their velocities, and are not the gradient of any scalar potential. In 3 dimensions, this implies that they have a non-vanishing curl, $\nabla \times \mathbf{F}(\mathbf{x})\neq0$. Following Berry and Shukla, these are referred to as \emph{curl-forces} \cite{berry2012classical}. 

Such curl-forces are usually associated with external forcing, which typically drives a thermal system out of equilibrium. These forces are not merely theoretical constructs: they can be generated in the lab, for example through an optical force field \cite{curtis2003structure, roichman2008influence,roichman2008optical, sun2009brownian, wu2009direct, liberal2013near, sukhov2017non, volpe2023roadmap}, or using feedback traps \cite{jun2012virtual}. 
A series of recent papers studied various properties of classical and quantum systems with curl-forces 
\cite{berry2012classical, berry2015hamiltonian, berry2016curl, berry2023quantum, berry2024quantising, berry2025six}. 
In particular, it was shown that certain two-dimensional systems with curl-forces can be described using Hamiltonian formalism. In this description, however, the Hamiltonian function differs from the canonical structure of kinetic plus potential energies \cite{berry2015hamiltonian}. The fact that some curl-forces can be described by a (non-standard) Hamiltonian raises a natural question: can the NESS of such systems be described by the Boltzmann distribution corresponding to their Hamiltonian? 

In this work, we show that unfortunately this is not the case: the Hamiltonian corresponding to a curl-force field cannot be used to easily find the steady state of systems driven out of equilibrium. However, we demonstrate how curl-forces can be used to give a clear qualitative description for currents arising in other, more obscure settings.
We begin in Sec.~\ref{sec:HCF} by a short introduction on Hamiltonian curl-forces. In Sec.~\ref{sec:StochDyn} we discuss the statistical mechanics of such systems when coupled to a heat bath, and explain why they do not relax to a Boltzmann distribution despite their Hamiltonian structure. We then consider the Klein-Kramers equation describing the stochastic evolution of the system.
In Sec.~\ref{sec:Transformation}, we present our main result: using a specific canonical transformation, we show that there is a mapping between thermal systems with Hamiltonian curl-forces and systems without any curl-forces that are instead simultaneously coupled to several heat baths at different temperatures, which is a different class of nonequilibrium systems. We study two examples of such systems in Sec.~\ref{Sec:Applications}, both featuring rotational currents solely due to the coupling to two separate heat baths, without any curl-forces. By analyzing $\nabla \times \mathbf{F}$ in the transformed, curl-force description of the same systems, we are able to provide a simple mechanical explanation for the origin of these currents. Finally, in Sec.~\ref{sec:discussion} we summarize our results and provide an outlook with open questions and possible generalizations.
